% ============================================================
% chapter3.tex  —  Chapitre 3 : Conception et Modélisation
% ============================================================

\chapter{Conception et Modélisation}


\noindent
The previous chapter described \emph{how} the platform was built.
This chapter describes \emph{what} was designed and \emph{why}
each component takes the form it does.
It presents the conceptual models, the database schema,
the API structure, and the interaction patterns
that define the DecisioBI platform.
The goal is to provide a complete, verifiable blueprint
of the system — one from which a reader could reconstruct
the implementation or evaluate its fitness for purpose.


% ============================================================
\section{Exigences Fonctionnelles}
% ============================================================

\noindent
Based on the gap analysis conducted in Chapter~1
and the target user profile (SME managers and analysts
in Sub-Saharan Africa), the following functional
requirements were identified.
Each requirement is expressed as a user story
in the format \textit{As a [role], I want [action]
so that [benefit].}

% --- Tableau des exigences fonctionnelles ---
\begin{table}[h!]
\centering
\caption{Exigences fonctionnelles principales.}
\label{tab:exigences}
\small
\begin{tabular}{|c|p{4.5cm}|p{4.5cm}|}
\hline
\textbf{ID} & \textbf{Exigence} & \textbf{Justification} \\
\hline
EF01 & As an admin, I want to register my organisation
      and invite users so that my team can collaborate. & Première utilisation — besoin d'espace de travail partagé. \\
\hline
EF02 & As an analyst, I want to upload CSV, Excel, and JSON files
      so that I can consolidate data from multiple sources. & Les PME utilisent principalement des fichiers plats. \\
\hline
EF03 & As an analyst, I want to preview data before import
      so that I can catch errors early. & Éviter les imports erronés. \\
\hline
EF04 & As an analyst, I want the system to automatically clean
      my data so that I don't spend hours on manual corrections. & Le nettoyage manuel est la principale source de retard. \\
\hline
EF05 & As an analyst, I want to review and validate cleaning results
      so that I retain control over my data. & Transparence et traçabilité des transformations. \\
\hline
EF06 & As a manager, I want to see conflicts between data sources
      so that I can trust the consolidated picture. & Conflits invisibles mènent à de mauvaises décisions. \\
\hline
EF07 & As a manager, I want to view KPIs automatically calculated
      from my data so that I can monitor performance in real time. & Remplacer les tableaux Excel manuels. \\
\hline
EF08 & As a manager, I want to receive alerts when a KPI
      crosses a threshold so that I can react quickly. & Détection proactive des anomalies. \\
\hline
EF09 & As a manager, I want to interact with data in natural language
      so that I don't need to learn SQL or BI tools. & Accessibilité pour les utilisateurs non techniques. \\
\hline
EF10 & As an admin, I want to control who can access what
      so that sensitive data remains protected. & Sécurité et conformité. \\
\hline
\end{tabular}
\end{table}


% ============================================================
\section{Exigences Non Fonctionnelles}
% ============================================================

\noindent
Beyond functional behaviour, the platform must satisfy
several non-functional requirements:

\begin{itemize}
  \item \textbf{Performance:} API response times under 200\,ms
        for standard operations; KPI calculation
        completes within seconds for datasets
        up to 100\,000 rows.
  \item \textbf{Scalability:} the architecture must support
        concurrent users from multiple organisations
        without degradation.
  \item \textbf{Usability:} the interface must be navigable
        by users with limited technical training;
        all labels are in French.
  \item \textbf{Security:} JWT-based authentication with
        role-based access control;
        data isolation between organisations;
        password hashing and account lockout.
  \item \textbf{Reliability:} background tasks (ingestion,
        cleaning) must not crash the API;
        failed tasks must be retryable.
  \item \textbf{Accessibility:} the frontend must function
        under low-bandwidth conditions
        typical of West African networks.
\end{itemize}


% ============================================================
\section{Modélisation des Cas d'Utilisation}
% ============================================================

\noindent
The platform involves three principal actor categories
and a machine actor (the AI engine).
Table~\ref{tab:acteurs} defines these actors.

\begin{table}[h!]
\centering
\caption{Acteurs du système.}
\label{tab:acteurs}
\begin{tabular}{|l|p{8cm}|}
\hline
\textbf{Acteur} & \textbf{Description} \\
\hline
Administrateur & Gère l'organisation, les utilisateurs,
                 les rôles et la configuration système.
                 Seul acteur pouvant créer des comptes
                 et modifier les permissions. \\
\hline
Analyste & Importe, nettoie et transforme les données.
           Peut créer des KPI et des règles de nettoyage.
           Ne peut pas gérer les utilisateurs. \\
\hline
Manager & Consulte les tableaux de bord, les KPI,
          les alertes et les analyses IA.
          Accès en lecture seule aux données. \\
\hline
Moteur IA & Génère des interprétations de KPI,
            détecte les anomalies et répond
            aux questions en langage naturel.
            Acteur non humain, déclenché par le système. \\
\hline
\end{tabular}
\end{table}

\noindent
The principal use cases are organised by module
in Table~\ref{tab:cas_utilisation}.

\begin{table}[h!]
\centering
\caption{Principaux cas d'utilisation.}
\label{tab:cas_utilisation}
\small
\begin{tabular}{|c|l|l|l|}
\hline
\textbf{ID} & \textbf{Cas d'utilisation} & \textbf{Acteur principal} & \textbf{Module} \\
\hline
CU01 & S'inscrire / Créer l'organisation & Admin & M9 \\
CU02 & Se connecter & Tous & M9 \\
CU03 & Gérer les utilisateurs & Admin & M9 \\
\hline
CU04 & Importer un fichier (CSV/Excel/JSON) & Analyste & M1 \\
CU05 & Aperçu des données avant import & Analyste & M1 \\
CU06 & Appliquer un template ERP & Analyste & M1 \\
\hline
CU07 & Exécuter un pipeline de nettoyage & Analyste & M2 \\
CU08 & Valider les résultats de nettoyage & Analyste & M2 \\
CU09 & Voir la comparaison avant/après & Manager & M2 \\
\hline
CU10 & Détecter les conflits inter-sources & Système & M3 \\
CU11 & Résoudre un conflit & Analyste & M3 \\
CU12 & Consulter le journal d'audit & Admin & M3 \\
\hline
CU13 & Consulter un KPI & Manager & M4 \\
CU14 & Créer un KPI personnalisé & Analyste & M4 \\
CU15 & Configurer une alerte KPI & Analyste & M4 \\
\hline
CU16 & Construire un tableau de bord & Manager & M5 \\
CU17 & Sauvegarder une vue personnalisée & Manager & M5 \\
\hline
CU18 & Demander une analyse IA & Manager & M6 \\
CU19 & Exécuter la détection d'anomalies & Système & M7 \\
CU20 & Consulter les anomalies détectées & Manager & M7 \\
\hline
CU21 & Poser une question au chatbot & Manager & M8 \\
CU22 & Démarrer une session de conversation & Manager & M8 \\
\hline
\end{tabular}
\end{table}


% ============================================================
\section{Modèle Conceptuel de Données}
% ============================================================

\noindent
The database design follows a relational model
with 27 tables distributed across 10 modules.
The schema is presented below at the conceptual level,
focusing on entities, attributes, and relationships.

% --- 3.4.1 Authentification et Organisation ---
\subsection{Module d'Authentification (M9)}

\noindent
The authentication module is built around four entities
that implement a multi-tenant RBAC model:

\begin{itemize}
  \item \textbf{Organization} — represents a company
        or workspace. Each registration creates
        one organisation. Fields: name, sector, size,
        country, created\_by, timestamps.
  \item \textbf{UserProfile} — extends Django's built-in
        User with role (admin/analyst/viewer),
        department, phone, MFA settings,
        failed login tracking, and account lockout
        (5 attempts $\rightarrow$ 15-minute lockout).
        Linked to Organization via foreign key.
  \item \textbf{AuthToken} — API tokens for programmatic
        access. Supports scopes, IP whitelisting,
        rate limiting (1\,000 requests default),
        and usage tracking.
  \item \textbf{Notification} — in-app notifications
        for pipeline events (file loaded, cleaning
        in progress, cleaning completed).
\end{itemize}

% --- 3.4.2 Ingestion ---
\subsection{Module d'Ingestion (M1)}

\noindent
Two entities manage the data acquisition process:

\begin{itemize}
  \item \textbf{DataSource} — metadata record for each
        uploaded file. Tracks source type (CSV, Excel, JSON),
        file size, row/column counts, encoding,
        checksum (MD5), status (pending $\rightarrow$
        processing $\rightarrow$ completed $\rightarrow$ failed),
        schema version, and lineage information
        (parent source for versioning).
  \item \textbf{RawData} — individual rows stored as
        JSON documents. Each row is linked to its
        DataSource and identified by row number.
        A unique constraint on (source, row\_number)
        prevents duplicate storage.
\end{itemize}

\noindent
An \textbf{IngestionJob} entity tracks the state
of asynchronous ingestion tasks, including
Celery task ID, progress percentage, and error messages.

% --- 3.4.3 Nettoyage ---
\subsection{Module de Nettoyage (M2)}

\noindent
The cleaning module introduces three entities:

\begin{itemize}
  \item \textbf{CleaningRule} — a reusable cleaning operation
        defined by a rule type (18 available types),
        column targeting (pattern or explicit list),
        parameters, and priority. Rules track
        execution count and success rate.
  \item \textbf{CleaningPipeline} — an ordered collection
        of rules applied as a sequence.
        Supports quality gates (minimum quality score,
        maximum missing rate) and source-type scoping.
  \item \textbf{CleaningJob} — execution record
        linking a source to a rule or pipeline.
        Tracks rows processed, affected, skipped,
        and failed; duration; memory usage;
        and whether the job was auto-triggered.
  \item \textbf{CleanedData} — stores the cleaned version
        of each row, linked to the original RawData.
        Includes change tracking (JSON list of modifications),
        quality score, and validation status.
\end{itemize}

% --- 3.4.4 Conflits ---
\subsection{Module de Résolution des Conflits (M3)}

\noindent
Three entities manage inter-source conflicts:

\begin{itemize}
  \item \textbf{ConflictType} — lookup table defining
        six conflict categories: duplicate records,
        missing values, data type issues,
        format inconsistencies, cross-source contradictions,
        and custom types. Each type specifies
        severity, auto-resolve capability,
        and default resolution strategy.
  \item \textbf{Conflict} — represents a detected conflict
        instance. Links to a DataSource and ConflictType.
        Tracks affected columns, row IDs,
        status (detected $\rightarrow$ investigating
        $\rightarrow$ resolving $\rightarrow$ resolved),
        priority, acknowledgement, assignment,
        due date, and impact score.
  \item \textbf{ConflictResolution} — audit record
        of how a conflict was resolved.
        Stores the resolution method (7 options:
        manual override, auto merge, default value,
        user selected, majority vote, latest value,
        discard), chosen value, alternative values,
        confidence score, and approval chain.
\end{itemize}

\noindent
An \textbf{ActivityLog} entity provides a comprehensive
audit trail of all user actions across the platform,
capturing user identity, action type, resource,
IP address, user agent, response time,
and risk score.

% --- 3.4.5 KPI ---
\subsection{Module KPI (M4/M5)}

\noindent
Three entities form the analytical core:

\begin{itemize}
  \item \textbf{KPI} — defines a key performance indicator
        with formula (SQL, Python, or Excel-style),
        target value, thresholds (warning, critical),
        frequency (daily to yearly), category,
        source table, measure and dimension columns,
        and optional parent-child hierarchy.
  \item \textbf{KPICalculation} — stores a computed value
        for a specific time period.
        Tracks previous value, variance (absolute
        and percentage), status (on target, warning,
        critical), breakdown details, data quality score,
        execution time, and optional forecast value
        with confidence interval.
  \item \textbf{KPIAlert} — configures threshold-based
        or anomaly-based alerts.
        Supports multi-channel notification
        (email, webhook, Slack), cooldown periods,
        escalation policies, and acknowledgement tracking.
\end{itemize}

% --- 3.4.6 Tableau de bord ---
\subsection{Module Tableau de bord (M5)}

\noindent
Four entities manage dashboard configuration:

\begin{itemize}
  \item \textbf{Dashboard} — defines a dashboard layout
        with grid columns (default 12),
        refresh interval, default filters,
        access control (public, allowed roles, allowed users),
        and view count tracking.
  \item \textbf{Widget} — individual dashboard component
        with type (metric card, line chart, bar chart,
        pie chart, table, gauge, scatter plot, heatmap),
        position (x, y), dimensions (width, height),
        data source reference, and drill-down capability.
  \item \textbf{Visualization} — reusable chart template
        with axis configuration, grouping, filtering,
        aggregation methods, style configuration,
        and usage tracking.
  \item \textbf{PreferenceUtilisateur} — per-user settings:
        visible columns, KPI order, dashboard layout,
        default period, currency (FCFA), number format
        (French/English).
  \item \textbf{VuePersonnalisee} — user-saved dashboard
        configurations (maximum 20 per user),
        with default and shared flags.
\end{itemize}

% --- 3.4.7 Intelligence artificielle ---
\subsection{Modules d'Intelligence Artificielle (M6/M7/M8)}

\noindent
Three entities manage AI-powered features:

\begin{itemize}
  \item \textbf{AIAnalysis} — records an AI analysis session:
        model provider, model name, prompt, context data,
        response, token usage, processing time,
        cost, and status.
  \item \textbf{AIInsight} — individual findings extracted
        from an analysis: type (trend, anomaly, pattern,
        correlation, recommendation, warning, opportunity),
        severity, supporting data, statistical significance,
        recommended action, estimated impact, and urgency.
  \item \textbf{AnomalyModel} / \textbf{Anomaly} /
        \textbf{AnomalyAlert} — ML model registry
        and detected anomaly records with contribution scores,
        severity classification, and multi-channel alerts.
\end{itemize}

% --- 3.4.8 Chatbot ---
\subsection{Module Chatbot (M8)}

\noindent
Three entities support conversational interaction:

\begin{itemize}
  \item \textbf{ChatSession} — represents a conversation
        with a title and status (active/archived).
  \item \textbf{ChatMessage} — individual messages
        within a session, with role (user/assistant),
        content, and intent metadata.
  \item \textbf{QueryHistory} — analytics on user queries:
        intent, entities, sentiment, execution time,
        and satisfaction score.
\end{itemize}


% ============================================================
\section{Schéma de la Base de Données — Diagramme Entité-Relation}
% ============================================================

\noindent
Figure~\ref{fig:erd} presents a simplified
entity-relationship diagram of the platform.
The 27 tables are organised into five clusters
围绕 the central DataSource entity.

\begin{figure}[h!]
\centering
\fbox{\parbox{0.92\textwidth}{\centering
\vspace{0.5cm}
\textbf{Diagramme Entité-Relation — Vue synthétique}

\vspace{0.3cm}
\small

\textsc{Authentification}
\begin{tabular}{|l|l|l|}
\hline
Organization & $\leftarrow$ created\_by -- User & \\
UserProfile & $\leftrightarrow$ User (1:1) & \\
UserProfile & $\rightarrow$ Organization (N:1) & \\
AuthToken & $\rightarrow$ User (N:1) & \\
\hline
\end{tabular}

\vspace{0.3cm}
\textsc{Pipeline de Données}

\begin{tabular}{|l|l|l|}
\hline
DataSource & $\leftarrow$ uploaded\_by -- User & \\
RawData & $\rightarrow$ DataSource (N:1) & \\
CleaningJob & $\rightarrow$ DataSource (N:1) & \\
CleaningJob & $\rightarrow$ CleaningRule (N:1) & \\
CleanedData & $\rightarrow$ CleaningJob (N:1) & \\
CleanedData & $\rightarrow$ RawData (N:1) & \\
\hline
\end{tabular}

\vspace{0.3cm}
\textsc{Conflits \& Résolution}

\begin{tabular}{|l|l|l|}
\hline
Conflict & $\rightarrow$ DataSource (N:1) & \\
Conflict & $\rightarrow$ ConflictType (N:1) & \\
ConflictResolution & $\rightarrow$ Conflict (N:1) & \\
\hline
\end{tabular}

\vspace{0.3cm}
\textsc{Analyse}

\begin{tabular}{|l|l|l|}
\hline
KPI & $\leftarrow$ owner -- User & \\
KPICalculation & $\rightarrow$ KPI (N:1) & \\
KPIAlert & $\rightarrow$ KPI (N:1) & \\
Widget & $\rightarrow$ Dashboard (N:1) & \\
Dashboard & $\leftarrow$ created\_by -- User & \\
AIAnalysis & $\rightarrow$ DataSource (N:1) & \\
AIAnalysis & $\rightarrow$ KPI (N:1) & \\
AIInsight & $\rightarrow$ AIAnalysis (N:1) & \\
Anomaly & $\rightarrow$ AnomalyModel (N:1) & \\
\hline
\end{tabular}

\vspace{0.3cm}
\textsc{Chatbot}

\begin{tabular}{|l|l|l|}
\hline
ChatMessage & $\rightarrow$ ChatSession (N:1) & \\
QueryHistory & $\rightarrow$ ChatSession (N:1) & \\
\hline
\end{tabular}

\vspace{0.5cm}
}}
\caption{Diagramme Entité-Relation simplifié (27 tables).}
\label{fig:erd}
\end{figure}


% ============================================================
\section{Modélisation des Classes}
% ============================================================

\noindent
At the implementation level, the backend follows
a consistent class structure per module.
Figure~\ref{fig:classes} presents the principal classes
and their relationships.

\begin{figure}[h!]
\centering
\fbox{\parbox{0.92\textwidth}{\centering
\vspace{0.5cm}
\textbf{Diagramme de Classes — Structure orientée objet}

\vspace{0.3cm}
\small

\texttt{[Service Layer]}

\begin{tabular}{|l|l|}
\hline
\texttt{KPICalculationService} & Calcule les KPI (SQL/Python/Excel) \\
\texttt{KPIAnomalyDetectionService} & Détecte les anomalies (Z-score, Isolation Forest, IQR) \\
\texttt{KPIForecastingService} & Prévisions par régression linéaire \\
\texttt{KPIAlertingService} & Envoie les alertes (email, webhook, Slack) \\
\texttt{FilterService} & Aliasage bilingue des champs (FR/EN) \\
\texttt{PivotService} & Tables croisées dynamiques \\
\texttt{NotificationService} & Notifications en 5 étapes \\
\hline
\end{tabular}

\vspace{0.3cm}
\texttt{[Nettoyage — Services de nettoyage]}

\begin{tabular}{|l|l|}
\hline
\texttt{StructuralCleaner} & Nettoyage structurel (colonnes, lignes) \\
\texttt{DateCleaner} & Parsing et normalisation des dates \\
\texttt{MontantCleaner} & Normalisation des montants \\
\texttt{TextCleaner} & Nettoyage du texte (standardisation) \\
\texttt{CoherenceChecker} & Vérification de cohérence \\
\texttt{ContextChecker} & Vérification contextuelle \\
\texttt{QualityScorer} & Calcul du score de qualité \\
\texttt{MLCleaner} & Profilage ML optionnel \\
\hline
\end{tabular}

\vspace{0.3cm}
\texttt{[API Layer] — Vues Django REST Framework}

\begin{tabular}{|l|l|}
\hline
\texttt{DataSourceUploadView} & Upload de fichiers \\
\texttt{CleaningApplyAsyncView} & Nettoyage asynchrone \\
\texttt{ConflictDetectionAPIView} & Détection de conflits \\
\texttt{KPIEngineAPIView} & Calcul de KPI \\
\texttt{PivotTableView} & Tables pivot \\
\texttt{InterpretKpisView} & Interprétation IA \\
\texttt{IsolationForestRunView} & Détection d'anomalies \\
\texttt{ChatSessionMessageListCreateView} & Chatbot \\
\hline
\end{tabular}

\vspace{0.5cm}
}}
\caption{Diagramme de classes — couches applicatives.}
\label{fig:classes}
\end{figure}


% ============================================================
\section{Modélisation des Interactions — Diagrammes de Séquence}
% ============================================================

\noindent
This section presents the sequence diagrams
for the three most critical workflows.

% --- 3.7.1 Flux d'import ---
\subsection{Flux d'Import de Données}

\noindent
The data import workflow involves five participants:
the User (analyst or admin), the Frontend,
the API, the Celery Worker, and the Database.

\begin{enumerate}
  \item The user selects a file and clicks \textit{Importer}.
  \item The frontend sends a POST request to
        \texttt{/api/ingestion/sources/async-upload/}
        with the file and metadata.
  \item The API validates the file format and size,
        creates a DataSource record with status
        \textit{pending}, and enqueues a Celery task.
  \item The API returns a job ID immediately
        (HTTP 202 Accepted).
  \item The Celery worker picks up the task,
        updates the DataSource status to
        \textit{processing}, parses the file,
        and inserts rows into the RawData table.
  \item On completion, the worker updates the status
        to \textit{completed} and fires a Django signal.
  \item The signal handler creates a notification
        (\textit{ingestion\_completed}) and,
        if auto-cleaning is enabled, enqueues
        a cleaning task.
  \item The frontend polls \texttt{/api/ingestion/jobs/<id>/}
        every 3 seconds and updates the UI
        with progress information.
\end{enumerate}

% --- 3.7.2 Flux de nettoyage ---
\subsection{Flux de Nettoyage Automatisé}

\noindent
The cleaning workflow is triggered automatically
after ingestion or manually by the user:

\begin{enumerate}
  \item A cleaning task is enqueued
        (either by signal after ingestion or by
        a POST to \texttt{/api/nettoyage/sources/<id>/apply-async/}).
  \item The Celery worker loads the raw data
        and the cleaning pipeline.
  \item Each rule in the pipeline is applied sequentially:
        structural cleaning, date parsing,
        amount normalisation, text standardisation,
        coherence checking.
  \item Quality scores are computed at each step.
  \item Results are stored in the CleanedData table
        with a before/after comparison.
  \item A cleaning job record tracks rows processed,
        affected, and skipped.
  \item The user reviews the comparison
        via \texttt{/api/nettoyage/jobs/<id>/comparison/}
        and either validates or rejects the results.
  \item Validation triggers a signal that enqueues
        conflict detection (M3).
\end{enumerate}

% --- 3.7.3 Flux de calcul de KPI ---
\subsection{Flux de Calcul de KPI}

\noindent
KPI calculation is event-driven:

\begin{enumerate}
  \item A data change event fires
        (after cleaning validation or conflict resolution).
  \item The event handler queries all active KPIs
        whose source matches the changed data.
  \item For each KPI, the KPICalculationService
        evaluates the formula:
        \begin{itemize}
          \item SQL formulas are executed
                against the database.
          \item Python formulas are evaluated
                in a sandboxed environment.
          \item Excel formulas are parsed
                and computed locally.
        \end{itemize}
  \item The calculated value is compared against
        the target, previous period, and thresholds.
  \item A KPICalculation record is created
        with variance analysis and status classification.
  \item If a threshold is breached,
        the KPIAlertingService is invoked:
        it checks cooldown periods,
        selects notification channels,
        and dispatches alerts.
  \item Anomaly detection runs concurrently:
        Z-score analysis flags values
        exceeding 2.5 standard deviations,
        and the Isolation Forest model
        identifies multivariate outliers.
  \item Results are pushed to the frontend
        via the dashboard API
        (\texttt{/api/kpi/dashboard/}).
\end{enumerate}


% ============================================================
\section{Conception de l'API REST}
% ============================================================

\noindent
The platform exposes a RESTful API
organised under the \texttt{/api/} prefix.
Table~\ref{tab:api} summarises the endpoint structure
by module.

\begin{table}[h!]
\centering
\caption{Structure de l'API REST.}
\label{tab:api}
\small
\begin{tabular}{|l|l|l|}
\hline
\textbf{Module} & \textbf{Préfixe} & \textbf{Endpoints principaux} \\
\hline
Auth (M9) & \texttt{/api/auth/} & register, login, logout, \\
& & token, profile, admin/users, \\
& & password-reset, notifications \\
\hline
Ingestion (M1) & \texttt{/api/ingestion/} & sources, upload, async-upload, \\
& & preview, templates, jobs, raw-data \\
\hline
Nettoyage (M2) & \texttt{/api/nettoyage/} & rules, pipelines, jobs, \\
& & suggestions, preview, apply, \\
& & apply-async, comparison, validate \\
\hline
Conflits (M3) & \texttt{/api/conflits/} & types, conflicts, resolutions, \\
& & activity-log, detect \\
\hline
KPI (M4/M5) & \texttt{/api/kpi/} & kpis, auto/detect, calculate\_now, \\
& & history, anomaly\_detection, \\
& & forecast, variance\_analysis, \\
& & alerts, workbench/pivot, \\
& & workbench/metric, filtres-disponibles \\
\hline
Dashboard (M5) & \texttt{/api/dashboard/} & dashboards, widgets, preferences, \\
& & vues, auto-build, add-widget, \\
& & analytics, preview \\
\hline
IA (M6) & \texttt{/api/ia/} & interpret-kpis, history, \\
& & sources/<id>/kpis \\
\hline
Anomalies (M7) & \texttt{/api/anomalies/} & detections, ml-models, \\
& & isolation\_forest/run \\
\hline
Chatbot (M8) & \texttt{/api/chatbot/} & sessions, messages \\
\hline
\end{tabular}
\end{table}

\noindent
All endpoints follow DRF conventions:
\texttt{GET} for listing/detail,
\texttt{POST} for creation,
\texttt{PUT/PATCH} for updates,
\texttt{DELETE} for removal.
Responses are paginated by default (20 items per page).
Errors return structured JSON with \texttt{detail}
and optional \texttt{errors} fields.


% ============================================================
\section{Modèle de Sécurité}
% ============================================================

\noindent
Security is designed at three layers,
as illustrated in Figure~\ref{fig:security}.

\begin{figure}[h!]
\centering
\fbox{\parbox{0.85\textwidth}{\centering
\vspace{0.5cm}
\textbf{Modèle de Sécurité — Trois couches}

\vspace{0.3cm}
\small

\textbf{Couche 1 : Authentification}
\begin{tabular}{|l|l|}
\hline
JWT Access Token & Valide 60 minutes \\
JWT Refresh Token & Valide 7 jours, rotation activée \\
API Token & SHA-256 hash, whiteliste IP, rate limit \\
Verrouillage & 5 tentatives échouées $\rightarrow$ 15 min \\
\hline
\end{tabular}

\vspace{0.3cm}
\textbf{Couche 2 : Autorisation (RBAC)}
\begin{tabular}{|l|c|c|c|}
\hline
Capacité & Admin & Analyste & Viewer \\
\hline
Lire le tableau de bord & $\checkmark$ & $\checkmark$ & $\checkmark$ \\
Lire les KPI & $\checkmark$ & $\checkmark$ & $\checkmark$ \\
Importer des données & $\checkmark$ & $\checkmark$ & $\times$ \\
Créer des KPI & $\checkmark$ & $\checkmark$ & $\times$ \\
Gérer les utilisateurs & $\checkmark$ & $\times$ & $\times$ \\
Utiliser le chatbot & $\checkmark$ & $\checkmark$ & $\checkmark$ \\
\hline
\end{tabular}

\vspace{0.3cm}
\textbf{Couche 3 : Isolation des Données}
\begin{tabular}{|l|}
\hline
Multi-tenant : chaque organisation ne voit que ses données \\
Requêtes SQL paramétrées (pas d'injection possible) \\
Sandbox Python pour les formules KPI \\
Hachage bcrypt des mots de passe \\
Hash SHA-256 des tokens API \\
\hline
\end{tabular}

\vspace{0.5cm}
}}
\caption{Modèle de sécurité à trois couches.}
\label{fig:security}
\end{figure}


% ============================================================
\section{Modèle de Déploiement}
% ============================================================

\noindent
The deployment architecture separates the frontend
and backend into independent scaling domains:

\begin{itemize}
  \item \textbf{Frontend:} static assets served by
        Cloudflare Workers (serverless edge),
        providing global low-latency access.
        Built with Vite, served as a single-page application.
  \item \textbf{Backend:} Django application server
        behind a reverse proxy,
        with PostgreSQL and Redis on the same machine
        or nearby network.
  \item \textbf{Workers:} Celery processes running
        as systemd services,
        handling ingestion, cleaning, conflict detection,
        KPI calculation, and anomaly detection.
  \item \textbf{Communication:} the frontend proxies
        \texttt{/api} requests to the backend
        in development; in production,
        API calls go directly to the backend's
        public URL via \texttt{VITE\_API\_BASE\_URL}.
\end{itemize}

\noindent
This architecture is deliberately simple.
It avoids container orchestration (Kubernetes),
service meshes, and other infrastructure complexity
that would be disproportionate to the project's scale.
A single server running Django, PostgreSQL, Redis,
and Celery workers is sufficient for the initial deployment,
with the option to scale horizontally
by adding worker processes or database replicas
if demand grows.


% ============================================================
\section{Résumé du Chapitre}
% ============================================================

\noindent
This chapter has presented the complete design
of the DecisioBI platform.
Ten functional requirements were derived from the gaps
identified in Chapter~1,
and six non-functional requirements
were established to ensure quality.
The data model comprises 27 tables
organised into 10 modules,
with a central DataSource entity
connecting ingestion, cleaning, conflict resolution,
KPI calculation, and AI analysis.
The API exposes 70+ endpoints
following REST conventions with JWT authentication
and role-based access control.
Three critical workflows — import, cleaning,
and KPI calculation — were modelled as sequence diagrams
to verify that the design supports
the required user interactions.
The deployment architecture is intentionally minimal,
prioritising simplicity and reliability
over infrastructure sophistication.

\vspace{0.4cm}
\noindent
The next chapter presents the implementation itself,
showing how these design decisions
translate into working code,
validated by a comprehensive test suite.

\vspace{0.7cm}
\begin{center}
  \textcolor{BITred!40}{\rule{0.35\linewidth}{0.5pt}}
  \hspace{0.5em}
  \textcolor{BITred}{$\star$}
  \hspace{0.5em}
  \textcolor{BITred!40}{\rule{0.35\linewidth}{0.5pt}}
\end{center}
